\chapter{Etude conceptuelle}

\section*{Modélisation}
\addcontentsline{toc}{section}{Modélisation}
L'approche conceptuelle adoptée repose sur la modélisation UML :
\begin{itemize}
	\item \textbf{Diagramme de classes} : structuration des entités principales (User, Admin, Enseignant, Étudiant, Cours, Demande, Notification, etc.) et de leurs relations ;
	\item \textbf{Cas d'utilisation} : identification des fonctionnalités par acteur ;
	\item \textbf{Diagrammes de séquence} : description des échanges entre contrôleurs, services et repositories.
\end{itemize}

Le backend suit une architecture 3-tiers : couche présentation API, couche métier, couche persistance.

\section*{Répartition des tâches conceptuelles (Lot A)}
\addcontentsline{toc}{section}{Répartition des tâches conceptuelles (Lot A)}
\begin{tabularx}{\textwidth}{lX}
\toprule
\textbf{Axe} & \textbf{Tâches Lot A} \\
\midrule
Données & Modélisation des entités JPA (environ 25 entités) et des relations inter-modules. \\
API & Définition des endpoints REST, des objets de réponse et des règles de validation. \\
Administration & Conception des workflows de traitement des demandes, absences, réclamations et notifications. \\
\bottomrule
\end{tabularx}
